\section*{Chapter \ref{ch:g8:fronts-and-weather-maps} --- Air Masses, Fronts and Weather Maps}

\begin{solution}{exo:g8:fronts-and-weather-maps:1}
(a) Maritime tropical: warm and damp. (b) Continental polar: cold and
dry. (c) Continental tropical: hot and dry. (d) Maritime polar: cold and
damp.
\end{solution}

\begin{solution}{exo:g8:fronts-and-weather-maps:2}
Blue line with triangles: a \omterm{def:g8:fronts-and-weather-maps:front}{cold front}. Red line with half circles: a
\omterm{def:g8:fronts-and-weather-maps:front}{warm front}. Line joining places with the same pressure: an \omterm{def:g8:fronts-and-weather-maps:isobar}{isobar}. H: the
centre of an \omterm{def:g8:fronts-and-weather-maps:depression}{anticyclone} (high pressure). L: the centre of a \omterm{def:g8:fronts-and-weather-maps:depression}{depression}
(low pressure).
\end{solution}

\begin{solution}{exo:g8:fronts-and-weather-maps:3}
(a) True. (b) False: they arrive before it, far ahead, high in the sky.
(c) False: close \omterm{def:g8:fronts-and-weather-maps:isobar}{isobars} mean a strong wind. (d) True.
\end{solution}

\begin{solution}{exo:g8:fronts-and-weather-maps:4}
(a) \qty{980}{hPa}. (b) \qty{1034}{hPa}. (c) \qty{4}{hPa}. (d) $1034 -
980 = \qty{54}{hPa}$.
\end{solution}

\begin{solution}{exo:g8:fronts-and-weather-maps:5}
Thin streaks of \omterm{def:g4:clouds-rain-snow:cloud-families}{cirrus}; a grey veil of altostratus; the \omterm{def:g4:clouds-rain-snow:cloud}{clouds} come lower
and thicker; steady \omterm{def:g4:clouds-rain-snow:precipitation}{rain}; warmer, damper air.
\end{solution}

\begin{solution}{exo:g8:fronts-and-weather-maps:6}
At a \omterm{def:g8:fronts-and-weather-maps:front}{cold front}, the dense cold air pushes under the warm air like a
\emph{wedge}: the warm air is lifted on a \emph{steep} slope, rises fast
and builds tall \omterm{def:g4:clouds-rain-snow:cloud}{clouds}. At a \omterm{def:g8:fronts-and-weather-maps:front}{warm front}, the warm air climbs slowly over
the cold air on a \emph{gentle slope}: it rises little by little over a
long distance, and its \omterm{def:g4:clouds-rain-snow:cloud}{clouds} spread out in layers.
\end{solution}

\begin{solution}{exo:g8:fronts-and-weather-maps:7}
(a) Near L: the \omterm{def:g8:fronts-and-weather-maps:isobar}{isobars} are closer together there, and the arrows are
longer. (b) From the east (the arrows point west, a little towards the
south-west). (c) From the east (the arrows point west). (d)
Anticlockwise.
\end{solution}

\begin{solution}{exo:g8:fronts-and-weather-maps:8}
(a) $1016 - 1003 = \qty{13}{hPa}$. (b) A \omterm{def:g8:fronts-and-weather-maps:depression}{depression}: the pressure keeps
falling. (c) \omterm{def:g4:clouds-rain-snow:cloud}{Clouds}, wind and \omterm{def:g4:clouds-rain-snow:precipitation}{rain}, probably with a \omterm{def:g8:fronts-and-weather-maps:front}{front}.
\end{solution}

\begin{solution}{exo:g8:fronts-and-weather-maps:9}
(a) cT, continental tropical. (b) Hot and dry. (c) Over the \omterm{def:g2:water-on-earth:ocean}{sea} it takes
up water vapour: it becomes damper, and the air near the water is cooled
a little by the \omterm{def:g2:water-on-earth:ocean}{sea}.
\end{solution}

\begin{solution}{exo:g8:fronts-and-weather-maps:10}
(a) $4 - 1 = 3$ days. (b) $3 \div 3 = 1$ day for each ten years. (c) It
needs $5 - 1 = 4$ more days: about 40 years.
\end{solution}

\begin{solution}{exo:g8:fronts-and-weather-maps:11}
(a) About \qty{1350}{km} (four and a half \omterm{def:g2:drawing-to-map:scale}{scale} bars). (b) \qty{4}{hPa}
over \qty{80}{km}: $4 \div 80 \times 100 = \qty{5}{hPa}$ per
\qty{100}{km}. (c) $4 \div 140 \times 100 \approx \qty{2.9}{hPa}$ per
\qty{100}{km}. (d) Near L, where the pressure changes faster; yes, the
arrows are longer near L.
\end{solution}

\begin{solution}{exo:g8:fronts-and-weather-maps:12}
(a) The \omterm{def:g8:fronts-and-weather-maps:front}{cold front}: $300 \div 50 = 6$ hours. The \omterm{def:g8:fronts-and-weather-maps:front}{warm front} moves at
\qty{25}{km} per hour: $200 \div 25 = 8$ hours. The \omterm{def:g8:fronts-and-weather-maps:front}{cold front} arrives
2 hours earlier. (b) First town: towering \omterm{def:g4:clouds-rain-snow:cloud}{clouds}, heavy showers, gusts,
then cooler air and a clearing sky. Second town: lowering layered \omterm{def:g4:clouds-rain-snow:cloud}{clouds}
and a long, steady \omterm{def:g4:clouds-rain-snow:precipitation}{rain}, then warmer, damper air. (c) The second town,
reached by warm air.
\end{solution}

\begin{solution}{pb:g8:fronts-and-weather-maps:1}
\textbf{1.} In the warm air, between the two \omterm{def:g8:fronts-and-weather-maps:front}{fronts}. The \omterm{def:g8:fronts-and-weather-maps:front}{warm front} has
already passed.

\textbf{2.} From the south-west.

\textbf{3.} Towards the east: the triangles of a \omterm{def:g8:fronts-and-weather-maps:front}{cold front} lie on the
side towards which it moves, and they point east, towards the town.

\textbf{4.} Towering \omterm{def:g4:clouds-rain-snow:cloud}{clouds}, heavy showers, perhaps thunder and \omterm{def:g4:clouds-rain-snow:precipitation}{hail},
gusts of wind; then cooler air, a wind from the west or north-west and a
clearing sky.

\textbf{5.} At \qty{180}{km} at 06:00 and at \qty{300}{km} at 09:00.

\textbf{6.} $300 - 180 = \qty{120}{km}$.

\textbf{7.} $120 \div 3 = \qty{40}{km}$ per hour.

\textbf{8.} The town is at \qty{480}{km}: $480 - 300 = \qty{180}{km}$.

\textbf{9.} $180 \div 40 = 4.5$ hours, that is 4 hours and 30 minutes.

\textbf{10.} 09:00 + 4~h~30~min = \textbf{13:30}.

\textbf{11.} No: the showers should end at about 14:30, before the
kick-off.

\textbf{12.} $180 \div 30 = 6$ hours: at 15:00, just at the kick-off.
Forecast: ``The \omterm{def:g8:fronts-and-weather-maps:front}{cold front} should reach the town at about \textbf{13:30},
with an hour of heavy showers and gusts, then clearing skies: the match
can most likely be played. If the \omterm{def:g8:fronts-and-weather-maps:front}{front} slows down, it could arrive at
15:00 and delay the match.''
\end{solution}
